% from ega4-21a-fr.tex
\oldpage[IV-4]{268}
  dans $X$), $\mathrm{N}_{f_*(\sh{O}_{X'})|\sh{O}_X}(s')$ est une section de $\sh{O}_X$ au-dessus de $U$ (cf. (\textbf{II}, 6.5.1)).
  (On rappelle que la condition (II) est vérifiée pour tout morphisme fini $f$ lorsque $X$ est un préschéma localement noethérien normal (\emph{loc. cit.}).)
\end{enumerate}

On sait alors (\textbf{II}, 6.5.5) que pour tout $\sh{O}_{X'}$-Module inversible $\sh{L}'$ on définit la norme $\sh{L}=\mathrm{N}_{X'/X}(\sh{L}')$ (que nous écrirons aussi $\mathrm{N}(\sh{L}')$), qui est un $\sh{O}_X$-Module inversible.
En outre, pour tout ouvert $U$ de $X$ et toute section régulière $s'\in\Gamma(f^{-1}(U),\sh{O}_{X'})$, la norme $\mathrm{N}_{X'/X}(s')=\mathrm{N}_{f_*(\sh{O}_{X'})|\sh{O}_X}(s')$ (que nous écrirons aussi $\mathrm{N}(s')$) est un élément régulier de $\Gamma(U,\sh{O}_X)$;
on est en effet ramené aussitôt au cas où $U=X$ est affine, et alors la conclusion résulte de \sref{IV.21.5.2} dans l'hypothèse (I);
d'autre part, dans l'hypothèse (II), le fait que $\sh{R}(X)$ soit un $\sh{O}_X$-Module plat entraîne que la section $s'\otimes1$ de $\Gamma(U,f_*(\sh{O}_{X'})\otimes_{\sh{O}_X}\sh{R}(X))$ est aussi régulière (\hyperref[0.6.1.2]{$\mathbf{0}_{\mathrm I}$, 6.1.2}), et la conclusion résulte encore de \sref{IV.21.5.2} appliqué à l'anneau $\Gamma(U,\sh{R}(X))$, compte tenu de la définition de la norme d'une section (\textbf{II}, 6.5.3).
Cela étant, soit $u'$ une section méromorphe de $\sh{L}'$ au-dessus de $X'$;
le morphisme $f$ étant affine, tout point $x\in X$ admet un voisinage ouvert $U$ tel que $u'|f^{-1}(U)$ soit de la forme $t'/s'$, où $t'\in\Gamma(f^{-1}(U),\sh{L}')$ et $s'$ est une section régulière dans $\Gamma(f^{-1}(U),\sh{O}_{X'})$;
l'élément $\mathrm{N}(t')/\mathrm{N}(s')$ (où $\mathrm{N}(t')$ est la section de $\sh{L}$ définie dans (\textbf{II}, 6.5.3)) est alors une section méromorphe de $\sh{L}$ au-dessus de $U$ en vertu de ce qui précède, et il résulte des propriétés multiplicatives de la norme (\textbf{II}, 6.5.3.1) que cette section ne dépend que de $u'|f^{-1}(U)$ et non de son écriture sous la forme $t'/s'$;
pour la même raison, lorsque $U$ varie, les sections méromorphes de $\sh{L}$ au-dessus de $U$ ainsi définies sont les restrictions d'une même section méromorphe de $\sh{L}$ au-dessus de $X$, que l'on appelle la \emph{norme} de $u'$ et que l'on note $\mathrm{N}_{X'/X}(u')$ (ou simplement $\mathrm{N}(u')$).
L'application ainsi définie
\[
\label{IV.21.5.3.1}
  \mathrm{N}_{X'/X}:\Gamma(X',\sh{M}_{X'}(\sh{L}'))\to\Gamma(X,\sh{M}_X(\mathrm{N}_{X'/X}(\sh{L}')))
  \tag{21.5.3.1}
\]
prolonge la norme définie dans (\textbf{II}, 6.5.5.3);
si $u'$ est une section méromorphe régulière de $\sh{L}'$ au-dessus de $X'$, il résulte aussitôt de ce qui précède que $\mathrm{N}(u')$ est une section méromorphe régulière de $\sh{L}$ au-dessus de $X$, car (avec les mêmes notations) $\mathrm{N}(t')$ est régulière si $t'$ l'est.
Enfin, si $\sh{L}_1'$, $\sh{L}_2'$ sont deux $\sh{O}_{X'}$-Modules inversibles, $s_1'$ (resp. $s_2'$) une section méromorphe régulière de $\sh{L}_1'$ (resp. $\sh{L}_2'$) au-dessus de $X'$, on a, en vertu de ce qui précède et de (\textbf{II}, 6.5.3.1)
\[
\label{IV.21.5.3.2}
  \mathrm{N}(s_1'\otimes s_2')=\mathrm{N}(s_1')\otimes\mathrm{N}(s_2').
  \tag{21.5.3.2}
\]
\end{env}

\begin{env}[21.5.4]
\label{IV.21.5.4}
Supposons toujours que $f$ vérifie l'une des hypothèses I), II) de \sref{IV.21.5.3}.
Si $(\sh{L}_1',s_1')$, $(\sh{L}_2',s_2')$ sont deux couples formés chacun d'un $\sh{O}_{X'}$-Module inversible et d'une section méromorphe régulière de ce Module au-dessus de $X'$, qui en outre sont équivalents au sens de \sref{IV.21.2.10}, alors les couples $(\mathrm{N}(\sh{L}_1'),\mathrm{N}(s_1'))$ et $(\mathrm{N}(\sh{L}_2'),\mathrm{N}(s_2'))$ sont aussi équivalents, car un isomorphisme de $\sh{O}_{X'}$-Modules inversibles a pour norme un isomorphisme de leurs normes (\textbf{II}, 6.5.3), et on a vu plus haut que $\mathrm{N}_{X'/X}$ transforme les sections de $\Gamma(X',\sh{O}_{X'}^*)$ en celles de $\Gamma(X,\sh{O}_X^*)$.
On peut donc poser la définition suivante :
\end{env}

